\RequirePackage[l2tabu, orthodox]{nag}

%\documentclass[a4paper]{article}
%\documentclass[final, pdftex, a4paper, 12pt, openbib, ]{article}
\documentclass[final, pdftex, a4paper, openbib, ]{article}
\usepackage[utf8]{inputenc}
\usepackage[francais]{babel}
%\usepackage{lmodern}
\usepackage[T1]{fontenc}
%\usepackage{graphicx}
\usepackage{alltt}
\usepackage{float}
%\usepackage{times}
%\usepackage{a4wide}
\usepackage{upquote,textcomp}
\usepackage{geometry}
%\usepackage{hyperref}
\usepackage[
pdftex,
final,                      % if you do    want to have clickable-colorful links
pdfstartview = FitV,
linktocpage  = false,       % ToC, LoF, LoT place hyperlink on page number, rather than entry text
breaklinks   = true,        % so long urls are correctly broken across lines
% pagebackref  = false,     % add page number in bibliography and link to position in document where cited
]{hyperref}
\usepackage{mathtools}
\geometry{a4paper, portrait, margin=2cm}
%\addtolength{\textheight}{2 cm}
%\addtolength{\oddsidemargin}{-1cm}
%\addtolength{\topmargin}{-1cm}
%\addtolength{\textwidth}{1 cm}

%Good solution for monospaces
%\usepackage{pxfonts} % Or palatino or mathpazo, changes all fonts to something sans
%\usepackage{eulervm} % only changes math fonts, i checked
%\usepackage[ttdefault=true]{AnonymousPro} %Only changes tt fonts
% WHICH ONE TO CHOOSE?
\usepackage{pslatex}
%\usepackage[ttdefault=true]{AnonymousPro} %Only changes tt fonts
% \usepackage[varg, cmintegrals, cmbraces, ]{newtxtext,newtxmath} % libertine, uprightGreek (U.S.) or slantedGreek (ISO), 
% \usepackage{tgtermes}
% \usepackage{txfonts}
% \usepackage{mathptmx}
% \usepackage[scaled=.90]{helvet}
% \usepackage{courier}
% \usepackage{textcomp}     % required for special glyphs
% \usepackage{bm}   


\usepackage{caption}
\captionsetup[figure]{labelformat=empty}% redefines the caption setup of the figures environment in the beamer class.

%Inconsoloata, a little too light.
%\usepackage{inconsolata}
%\renewcommand{\ttdefault}{Consolas}
%\usepackage{fontspec} %Doesn't work with pdflatex
%\setmonofont{Consolas}



%\renewcommand*\familydefault{\ttdefault} %% Only if the base font of the document is to be typewriter style

%\newcommand\Fontvi{\fontsize{6}{7.2}\selectfont}

%Make tabularx center cells vertically
%\def\tabularxcolumn#1{m{#1}}
%\renewcommand{\tabularxcolumn}[1]{>{\small}m{#1}}

%%Syntax hilighting
%\usepackage{fancyvrb}
\usepackage{minted}
%\usepackage[newfloat]{minted}
\usemintedstyle{borland}
%\usepackage{etoolbox}
%\AtBeginEnvironment{minted}{\singlespacing%
%	\fontsize{14}{14}\selectfont}
%\newminted{java}{fontsize=\footnotesize}
\usepackage{caption}

%\usepackage{multicol}
%\usepackage{vwcol}
%\usepackage{lipsum}
%\usepackage{microtype}
%\setlength{\columnseprule}{0.4pt}
%\renewcommand{\columnseprulecolor}{\color{red}}
\newcommand{\BAD}[1]{{\color{red}#1}}
\newcommand{\GOOD}[1]{{\color{darkgreen}#1}}


\usepackage{graphicx}
\usepackage{colortbl,array}
\usepackage{tabularx}
\definecolor{warningbackground}{RGB}{252,226,158}

\newcommand{\alertwarningbox}[1]{
    \centering
    \begin{tabularx}{0.9\linewidth}{
        >{\columncolor{warningbackground}}c
        >{\columncolor{warningbackground}}X}
        \raisebox{\dimexpr2\baselineskip-\height}
        {\includegraphics[scale=0.8]{attention.pdf}}&
        \raisebox{\tabcolsep}{\strut}#1\raisebox{-\tabcolsep}{\strut}
    \end{tabularx}
}


\usepackage{xcolor}

%\definecolor{infobackground}{RGB}{217,237,247}
%\definecolor{infoforeground}{RGB}{58,135,173}
%\definecolor{infoborder}{RGB}{188,232,241}

\definecolor{infobackground}{RGB}{217,237,247}
\definecolor{infoforeground}{RGB}{30,50,70}
\definecolor{infoborder}{RGB}{30,50,70}

\definecolor{infobackground}{RGB}{217,237,247}
\definecolor{infoforeground}{RGB}{30, 80, 150}
\definecolor{infoborder}{RGB}{47, 87, 232}


\usepackage{environ}
\usepackage{tikz}
\usetikzlibrary{fit,backgrounds,calc}

\NewEnviron{alertinfo}[1]
{
\begin{center}
    \begin{tikzpicture}
    \node[inner sep=0pt,
          draw=infoborder,
          line width=1pt,
          fill=infobackground] (box) {\parbox[t]{0.99\textwidth}
        {%
            \begin{minipage}{.12\textwidth}
                \centering\tikz[scale=3]
                \node[scale=1]
                {
%                    \includegraphics[scale=0.25]{\images/information.pdf}
                    \includegraphics[scale=0.15]{\images/information.pdf}                    
                };
            \end{minipage}%
           \begin{minipage}{.86\textwidth}
%                \vskip 10pt
                \textbf{\textcolor{infoforeground}{\large #1}}\par\smallskip
                \textcolor{infoforeground}{\large \BODY}
                \par\smallskip
                \par\smallskip
            \end{minipage}\hfill
        }%
    };
    \end{tikzpicture}
\end{center}
}

\usepackage{titling}
\usepackage{soul}

\definecolor{Light}{gray}{.93}
\sethlcolor{Light}
\let\OldTexttt\texttt	
\renewcommand{\texttt}[1]{\OldTexttt{\hl{#1}}}

\newcommand{\codes}{../codes/TP8}
\newcommand{\images}{../images}

%\title{GBIAAL 4$^{\mbox{\`eme}}$ année \\ Devoir Surveillé --- Base de données \\ 13 janvier 2016  \\ 1 heure}
\title{IMA 3$^{\mbox{\`eme}}$ année\\ Programmation Avancé
	%\\TP1 Structures, Listes contiguës, Redirections}
}
\author{\huge \textbf{TP8 Arbres d'entiers}}
\setlength{\parindent}{0pt}
\pagestyle{empty}
\date{}


\begin{document}
%\vspace{-15cm}
\posttitle{\par\end{center}}
\setlength{\droptitle}{-45pt}
\maketitle
%\thispagestyle{empty}

%\section{Schéma conceptuel}
\vspace{-1.7cm}
\section{Objectifs}

\begin{itemize}
	\item Savoir créer et utiliser une bibliothèque d'arbres d'entiers
	\item Savoir coder quelques algorithmes de base sur les arbres
\end{itemize}

%\paragraph{Contexte et préparation : } Avant de commencer, récupérez le fichier \url{~wrudamet/public/IMA3/TP6/tp6_gdb.tgz} à la ligne de commande. Désarchivez\footnote{Vous pouvez utiliser la commande \texttt{tar} pour décompresser et désarchivez ce fichier} le dans votre répertoire de \textit{Programmation Avancé} (toujours à la ligne de commande !).

\paragraph{Préparation : } Avant de commencer, copier le répertoire \url{~wrudamet/public/IMA3/TP8/} qui contient :

\begin{itemize}
	\item Les fichiers \texttt{trees.c} et \texttt{trees.h}, contenant des fonctions de base de traitement d'arbres
	\item Un fichier \texttt{main.c} utilisant ces fonctions
	\item Un sous-répertoire \texttt{samples/} qui contient des fichiers de données \texttt{.txt} et des fichiers \texttt{.dot} et \texttt{.pdf} qui permettent de visualiser les arbres correspondant à ces données (section 2).
	\item Un sous répertoire \texttt{part2/} contenant un autre \texttt{main.c}.
\end{itemize}

%\begin{alertinfo}{}
%	\textbf{IMPORTANT : Ce TP sera réutilisé dans le futur. Il servira d'inspiration pour le TP8 et votre librairie sera entièrement repris pour le TP9}
%\end{alertinfo}

\section{Fonctions de base dans les ABR}
Dans cette première partie, il vous est demandé d'implémenter les fonctions suivantes non implémentées dans \texttt{tree.h}, en utilisant le \texttt{main} fourni pour tester :

\begin{enumerate}
	\item Ajout d'un élément dans un ABR (arbre binaire de recherche, donc ordonné !). (Fait en cours)
	\item Construction d'un ABR à partir d'un fichier de données. Pour cette fonction, on s'inspirera de la fonction \texttt{charge\_fichier} donnée dans le TP7 (attention, on ajoute des entiers et non des chaînes)
	\item Impression d'un ABR par parcours récursif (fait en cours)
	\item Vérifier que la fonction de parcours précédente donne la même suite de valeurs pour les trois fichiers \texttt{unspecified.txt}, \texttt{balanced.txt}, \texttt{degenerated.txt} alors que ce sont des arbres différents (voir les dessins dans le répertoire \texttt{samples/})
	\item Écrire une autre fonction d'impression qui imprime les couples (père, fils) d'un arbre donné. On fera attention à ne pas imprimer les sous-arbres vides
	\item Vérifier que l'impression donne des arbres correspondant aux fichiers pdf fournis dans le répertoire \texttt{samples/}
	\item Archiver le fichier \texttt{trees.o} dans une bibliothèque \texttt{libtrees.a} (commande \texttt{ar}).
\end{enumerate}


\section{ Impression des arbres sous forme graphique}
Les fichiers \texttt{.pdf} ont été générés à partir des fichiers \texttt{.dot} (fournis dans \texttt{samples/}).
Le format \texttt{.dot} est un format de représentation textuelle simple d'arbres\footnote{Et plus généralement de graphes exploitables par des logiciels tels que \url{http://www.graphviz.org}}.
Visualisez le format de ces fichiers dans votre éditeur de texte préféré ou par la commande \texttt{less}\footnote{digraph signifie \textit{directed graph}, les arbres en font partie du fait de leur orientation \texttt{pere->fils}}.
À partir de ce format \texttt{.dot}, la commande \texttt{dot} permet de générer une version visualisable (.pdf, .png, \ldots) comme suit (en pdf par exemple) :

\begin{minted}[mathescape=true,escapeinside=||,tabsize=4,
%fontsize=\footnotesize,
]{bash}
				dot -Tpdf samples/balanced.dot -o samples/balanced.pdf
\end{minted}

L'objectif est alors de générer de tels fichiers pour des arbres construits à partir de fichiers de données (tels que ceux fournis dans \texttt{samples/*.txt}).
					
\paragraph{Préparation :} Vous pouvez copier votre bibliothèque \texttt{libtrees.a} dans le répertoire \texttt{part2/}. Dans ce répertoire, un \texttt{main.c} plus conséquent est fourni qui :
\begin{enumerate}
	\item Génère, par appel à la fonction \texttt{generate\_dot}, les fichiers \texttt{.dot} correspondants aux fichiers de données dont les noms ont été passés en paramètre.
	\item Transforme ces fichiers \texttt{.dot} en \texttt{.pdf} par appel système à la commande \texttt{dot}.
\end{enumerate}

\paragraph{Travail à faire :}
\begin{enumerate}
	\item Programmer la fonction \texttt{recursive\_dot} qui traite le cas général
	\item Tester ensuite par la commande : \texttt{./trees ../samples/*.txt}
Vous devez retrouver les fichiers \texttt{.dot} et \texttt{.pdf} fournis.
\end{enumerate}

%\section{Annexes}
%\vspace{-0.5cm}

\end{document}